\section{Development Journey: Mistakes, Fixes, and Validation}
\label{sec:journey}

The methodological core of this project is not a single algorithm but a disciplined sequence of
corrections, each documented as \emph{mistake $\rightarrow$ how it was found $\rightarrow$ the fix
$\rightarrow$ how the fix was validated}. We summarize the sixteen entries below; the validation
philosophy throughout is that, absent ground truth, ``accurate'' means ``healthy lands in the published
range and discriminates pathology, with documented limits,'' and that uncertain norms/methods are
grounded in primary sources via adversarial literature search.

\paragraph{J1 --- A phantom norm debunked.} The in-code healthy \hahp{} of $0.97\pm0.02$ was unsourced; a
literature search traced it to an osteoporotic, all-female cohort. Measuring 12 healthy necks gave
$0.94\pm0.13$, within the true healthy range~\cite{tan2004,lee2012,kaur2025}. Recalibration dropped the
healthy false-positive rate from 17\% to 0\%.

\paragraph{J2 --- Body isolation by canal-cut.} Erosion/connected-component body isolation left a wedge
and a stuck ratio; the fix was to define the body as everything anterior to the canal.

\paragraph{J3 --- Endplate-line heights.} Image-axis heights failed on tilted, tapered cervical bodies;
PCA orientation $+$ endplate-line fitting collapsed the spread and is robust to tilt and taper.

\paragraph{J4 --- Honest scoping of the fracture screen.} Validated against RSNA-2022, the wedge metric
had near-zero power on non-compression fractures, so 5.2 was scoped to a vertebral-body compression
screen, with the RSNA negative result banked.

\paragraph{J5 --- Adopting SCIseg for myelomalacia.} Hand-rolled intensity thresholds could not separate
lesion from normal variation; we adopted the published SCIseg model~\cite{sciseg2024} and reframed our
job as the specificity check.

\paragraph{J6 --- An over-claim, corrected, and the keystone.} We initially mislabelled teammate
measurements as ``inaccurate'' by treating quality flags as clinical ones; correcting that, and running
the canonical code, localized the real problem to unstable corner-extrema landmarking (heights backwards,
Cobb kyphotic, slip over-firing) --- one keystone, not five bugs.

\paragraph{J7--J9 --- The corner fix, grounded and implemented.} An adversarial literature workflow
confirmed endplate-line fitting beats corner extrema~\cite{wang2023,chen2013}; implementing it fixed the
Cobb sign (lordotic) and the heights, while the C6/C7 endpoints remained noisy --- isolating
body-isolation quality as the next lever.

\paragraph{J10 --- Myelomalacia specificity measured.} SCIseg on the healthy cords gave 10/11 clean
($\approx$91\% specificity), and the lesion-to-level mapping was wired end-to-end.

\paragraph{J11 --- A negative result, kept.} Swapping the body isolation for the SPINEPS corpus border
made the Cobb \emph{worse}; we recorded the negative rather than tuning to hide it, and noted that without
GT we can only measure precision, not accuracy.

\paragraph{J12 --- The endplate-voxel breakthrough.} A research workflow found that SPINEPS already
predicts the endplate itself (instance label $100+X$); fitting the line to those voxels (method C1) beat
both corner-cut and corpus on every span and matched the literature lordosis mean, closing the C6/C7
precision gap (SD \SI{5.9}{\degree}). It also corrected the false supine-offset assumption.

\paragraph{J13 --- An integration fixture bug, caught by a test.} Integrating Group 5 into the
measurement service produced a failing test whose synthetic fixture had the anterior/posterior walls
swapped for the RAS convention; fixing the fixture (not the science) restored a green suite --- a reminder
that synthetic fixtures must match real anatomy.

\paragraph{J14--J15 --- Validation on real cohorts.} The pipeline was run on 12 healthy and 10
symptomatic necks; Group 3 separated at $p=0.0001$, Group 4-C1 directionally, and Group 1 was correctly
null on a non-compression cohort (Section~\ref{sec:results}).

\paragraph{J16 --- The harness as a bug detector.} The same validation caught that Group 2's DHI and bulge
read backwards and localized the cause to an over-measured denominator at the cervicothoracic junction ---
a real, sign-level bug a unit test on clean shapes would have missed.

\paragraph{J17--J21 --- Owning the group code; G1 recalibrations.} Consolidating ownership of the
measurement modules, four service-level defects were measured and fixed on real cohorts: an over-aggressive
vertebral-tilt cut (\SI{20}{\degree} flagged 88\% of healthy bodies; recalibrated to \SI{45}{\degree}), the
vertebral-height corner method (reading \hahp{} backwards, \(\sim\)1.08; replaced by the endplate-line
fit, \(\sim\)0.93), the disc bulge reference, and a debunked absolute disc-height cut (replaced by a cited
relative rule). A resolution-robustness check confirmed the mm metrics are stable from \SI{0.8}{} to
\SI{4}{\milli\meter} through-plane, while the canal-cut Cobb is not --- a third argument for the
endplate-voxel method.

\paragraph{J22--J23 --- The disc, scaled and stratified.} A 49-case within-MMCSD test (276 disc-levels,
lesion vs.\ non-lesion) replaced the confounded cross-dataset comparison. The disc \emph{signal} grade
proved a non-discriminator (AUC $0.50$), as did posterior bulge; the disc/VB AP-width ratio carried the
only real signal (AUC $0.62$). Critically, an apparent strong AP-width effect (AUC $0.79$) collapsed to
$0.61$ under level stratification --- the nuisance variable (lesions cluster at wider levels) had nearly
produced a false headline.

\paragraph{J24--J26 --- The alignment reversal.} Group~4 Cobb was the sharpest lesson. At $n{=}11$ it
showed a large directional effect ($d{=}0.76$); we resisted reporting it as validated and instead built a
representative healthy cohort (26 multi-site controls). The effect \emph{dissolved} ($d{=}0.28$, AUC
$0.57$, $p{=}0.32$): the original controls had been lordosis-biased, and supine positioning is a
cross-cohort confound. Alignment is a non-specific marker; the larger sample, not a larger easy-side
sample, settled it.

\paragraph{J25 --- End-to-end and demographics.} The disc group was wired into the orchestrator, patient
age/sex were threaded through to the interpretation layer (with a cited sex-specific dural-sac cut), and
the full pipeline was run end-to-end to produce structured reports cross-referenced against a deployed
front-end render.

\paragraph{Process lessons.} Several themes recur: (i) ground every clinical number in a primary source,
because plausible-looking in-code constants can trace to the wrong cohort; (ii) separate measurement
\emph{method} quality from segmentation/body-isolation quality --- they are different levers; (iii) keep
negative results; (iv) without ground truth, distinguish precision from accuracy and never tune to make a
method merely \emph{look} tighter; and (v) a validation harness that compares healthy to pathology is also
the cheapest, most honest bug detector available.
